\section*{Chapitre \ref{ch:iv1:speed-drills} --- Exercices de vitesse}
\markboth{\protect\omTocMark{}{Réponses~: Exercices de vitesse}}{\protect\omTocMark{}{Réponses~: Exercices de vitesse}}

\begin{solution}{iq:iv1:speed-drills:1}
\textbf{Feuille 1.} Corriger la feuille avec cette clé. Le score est le nombre de bonnes réponses moins deux fois le nombre de mauvaises.
\par\nopagebreak
% Generated by code/interviews-1/03-speed-drills/python/gen_drill.py. Do not edit. French copy by draft/pilot-quant-fr/ivgen_fr.py.
\begin{ivdrillkey}
\ivkey{1}{$1881$} & \ivkey{17}{$390$} & \ivkey{33}{$0.4$} & \ivkey{49}{$4.63$} & \ivkey{65}{$2760$} \\
\ivkey{2}{$750$} & \ivkey{18}{$75$} & \ivkey{34}{$0.27$} & \ivkey{50}{$2000$} & \ivkey{66}{$36$} \\
\ivkey{3}{$1071$} & \ivkey{19}{$20$} & \ivkey{35}{$3288$} & \ivkey{51}{$348$} & \ivkey{67}{$1675$} \\
\ivkey{4}{$570$} & \ivkey{20}{$1065$} & \ivkey{36}{$2590$} & \ivkey{52}{$\frac{5}{6}$}\rule[-9.5pt]{0pt}{23pt} & \ivkey{68}{$112$} \\
\ivkey{5}{$420$} & \ivkey{21}{$984$} & \ivkey{37}{$2499$} & \ivkey{53}{$1290$} & \ivkey{69}{$\frac{5}{3}$}\rule[-9.5pt]{0pt}{23pt} \\
\ivkey{6}{$26$} & \ivkey{22}{$\frac{13}{10}$}\rule[-9.5pt]{0pt}{23pt} & \ivkey{38}{$25$} & \ivkey{54}{$1241$} & \ivkey{70}{$195$} \\
\ivkey{7}{$23.86$} & \ivkey{23}{$43$} & \ivkey{39}{$20.96$} & \ivkey{55}{$3798$} & \ivkey{71}{$3168$} \\
\ivkey{8}{$2110$} & \ivkey{24}{$\frac{5}{12}$}\rule[-9.5pt]{0pt}{23pt} & \ivkey{40}{$12.4$} & \ivkey{56}{$224$} & \ivkey{72}{$893$} \\
\ivkey{9}{$55$} & \ivkey{25}{$\frac{2}{3}$}\rule[-9.5pt]{0pt}{23pt} & \ivkey{41}{$954$} & \ivkey{57}{$16$} & \ivkey{73}{$72$} \\
\ivkey{10}{$16$} & \ivkey{26}{$\frac{1}{8}$}\rule[-9.5pt]{0pt}{23pt} & \ivkey{42}{$5.8$} & \ivkey{58}{$540$} & \ivkey{74}{$1048$} \\
\ivkey{11}{$24.78$} & \ivkey{27}{$34$} & \ivkey{43}{$\frac{3}{25}$}\rule[-9.5pt]{0pt}{23pt} & \ivkey{59}{$1950$} & \ivkey{75}{$1053$} \\
\ivkey{12}{$\frac{2}{5}$}\rule[-9.5pt]{0pt}{23pt} & \ivkey{28}{$120$} & \ivkey{44}{$28$} & \ivkey{60}{$990$} & \ivkey{76}{$1448$} \\
\ivkey{13}{$627$} & \ivkey{29}{$473$} & \ivkey{45}{$1$} & \ivkey{61}{$\frac{5}{48}$}\rule[-9.5pt]{0pt}{23pt} & \ivkey{77}{$135$} \\
\ivkey{14}{$2191$} & \ivkey{30}{$204$} & \ivkey{46}{$42$} & \ivkey{62}{$1241$} & \ivkey{78}{$144$} \\
\ivkey{15}{$1355$} & \ivkey{31}{$1878$} & \ivkey{47}{$19$} & \ivkey{63}{$1211$} & \ivkey{79}{$0.04$} \\
\ivkey{16}{$3750$} & \ivkey{32}{$120$} & \ivkey{48}{$1280$} & \ivkey{64}{$2496$} & \ivkey{80}{$19$} \\
\end{ivdrillkey}


\textbf{Quand laisser une case vide.} Une réponse devinée, juste avec la probabilité $p$, ajoute en moyenne $p - 2(1 - p) = 3p - 2$ au score. Ce gain n'est positif que lorsque $p > \tfrac23$~: il faut donc laisser une question vide, sauf si l'on est assez sûr de soi.

\textbf{Le rythme.} Huit minutes pour 80 questions font 6 secondes par question. Passer une question dès que l'on ne voit pas de voie rapide, car une question lente coûte deux ou trois questions faciles.

\textbf{Où l'on perd des points.} Chaque question appartient à une famille. Repérer la famille de chaque question manquée ou passée, et s'entraîner sur cette seule famille.
% Generated by code/interviews-1/03-speed-drills/python/gen_drill.py. Do not edit. French copy by draft/pilot-quant-fr/ivgen_fr.py.
\ivfamily{Additionner et soustraire des nombres entiers}{3, 8, 15, 16, 20, 31, 35, 41, 54, 55, 62, 72, 74, 75}{Travailler depuis la gauche, d'abord les centaines, et tenir un total courant.}
\ivfamily{Additionner et soustraire des décimaux}{7, 11, 39, 40, 42, 49, 57, 79}{Aligner la virgule mentalement, ou multiplier par 100 puis diviser par 100.}
\ivfamily{Multiplier par un chiffre}{14, 21, 36, 48, 51, 60, 63, 65, 73, 76}{Décomposer le plus grand nombre, ou arrondir puis soustraire quand c'est plus court~: $389 \times 7 = 400 \times 7 - 11 \times 7 = 2723$.}
\ivfamily{Nombres avec un raccourci}{1, 2, 5, 13, 29, 37, 50, 53, 59, 64, 67, 71}{$\times 25$~: multiplier par 100, puis diviser par 4. $\times 15$~: ajouter la moitié et multiplier par 10. $\times 11$~: multiplier par 10 et ajouter le nombre lui-même. $\times 99$~: multiplier par 100 et soustraire le nombre lui-même. $(n - d)(n + d) = n^2 - d^2$.}
\ivfamily{Division exacte}{6, 10, 27, 38, 44, 47, 66, 68, 70, 80}{$\div 5$~: doubler puis diviser par 10. $\div 25$~: multiplier par 4 et diviser par 100. $\div 12$ et $\div 15$~: diviser d'abord par 3.}
\ivfamily{Fractions}{12, 22, 24, 25, 26, 43, 45, 52, 61, 69}{Utiliser le plus petit dénominateur commun, multiplier en croix, simplifier avant de multiplier. Les réponses sont sous forme irréductible.}
\ivfamily{Pourcentages}{4, 9, 17, 19, 28, 30, 32, 46, 56, 58}{Partir de 10~\% et monter~: $15\%$ vaut $10\% + 5\%$, et $12.5\%$ vaut un huitième. $p\%$ de $n$ égale $n\%$ de $p$.}
\ivfamily{Produits de décimaux}{18, 23, 33, 34, 77, 78}{Multiplier les chiffres et compter les décimales. Traiter $0.25$, $0.75$, $1.5$ et $2.5$ comme des fractions~: $\tfrac14$, $\tfrac34$, $\tfrac32$, $\tfrac52$.}


\iqwrong{Traiter les questions strictement dans l'ordre et refuser de passer. La pénalité sanctionne les réponses hâtives, et non les cases vides~: 60 bonnes réponses et aucune mauvaise rapportent davantage que 75 bonnes et 8 mauvaises, soit 60 contre 59.}
\iqfollow{Comment obtenir une nouvelle feuille~?}{On lance \texttt{gen\_drill.py} dans le dossier de code de ce chapitre avec une autre graine. Ses questions viennent des mêmes familles, et il écrit une nouvelle clé.}
\end{solution}

\begin{solution}{iq:iv1:speed-drills:2}
\textbf{Feuille 2.} Corriger la feuille avec cette clé. Le score est le nombre de bonnes réponses moins le nombre de mauvaises.
\par\nopagebreak
% Generated by code/interviews-1/03-speed-drills/python/gen_drill.py. Do not edit. French copy by draft/pilot-quant-fr/ivgen_fr.py.
\begin{ivdrillkey}
\ivkey{1}{$\frac{44}{35}$}\rule[-9.5pt]{0pt}{23pt} & \ivkey{11}{$8879$} & \ivkey{21}{$4225$} & \ivkey{31}{$9991$} & \ivkey{41}{$0.4375$} \\
\ivkey{2}{$13803$} & \ivkey{12}{$2025$} & \ivkey{22}{$\frac{27}{16}$}\rule[-9.5pt]{0pt}{23pt} & \ivkey{32}{$9936$} & \ivkey{42}{$2991$} \\
\ivkey{3}{$27.2$} & \ivkey{13}{$1925$} & \ivkey{23}{$\frac{7}{12}$}\rule[-9.5pt]{0pt}{23pt} & \ivkey{33}{$\frac{1}{3}$}\rule[-9.5pt]{0pt}{23pt} & \ivkey{43}{$93.5$} \\
\ivkey{4}{$85.5$} & \ivkey{14}{$\frac{37}{36}$}\rule[-9.5pt]{0pt}{23pt} & \ivkey{24}{$1829$} & \ivkey{34}{$14.4$} & \ivkey{44}{$0.45$} \\
\ivkey{5}{$0.128$} & \ivkey{15}{$5088$} & \ivkey{25}{$27.5$} & \ivkey{35}{$5988$} & \ivkey{45}{$4180$} \\
\ivkey{6}{$4376$} & \ivkey{16}{$3105$} & \ivkey{26}{$0.175$} & \ivkey{36}{$4489$} & \ivkey{46}{$0.3125$} \\
\ivkey{7}{$195.5$} & \ivkey{17}{$4439$} & \ivkey{27}{$31.5$} & \ivkey{37}{$3772$} & \ivkey{47}{$4649$} \\
\ivkey{8}{$0.375$} & \ivkey{18}{$2070$} & \ivkey{28}{$9135$} & \ivkey{38}{$\frac{65}{36}$}\rule[-9.5pt]{0pt}{23pt} & \ivkey{48}{$4096$} \\
\ivkey{9}{$9975$} & \ivkey{19}{$2736$} & \ivkey{29}{$4730$} & \ivkey{39}{$1517$} & \ivkey{49}{$\frac{10}{3}$}\rule[-9.5pt]{0pt}{23pt} \\
\ivkey{10}{$0.5625$} & \ivkey{20}{$2704$} & \ivkey{30}{$7474$} & \ivkey{40}{$1225$} & \ivkey{50}{$2030$} \\
\end{ivdrillkey}


\textbf{Quand laisser une case vide.} Avec $+1$ et $-1$, une réponse devinée, juste avec la probabilité $p$, ajoute en moyenne $p - (1 - p) = 2p - 1$ au score. Ce gain est positif lorsque $p > \tfrac12$~: la pénalité tolère donc plus de doute que sur la feuille 1, et l'on répond dès que l'on juge plus probable d'avoir raison que tort.

\textbf{Le rythme.} Six minutes pour 50 questions font 7.2 secondes par question, et beaucoup de questions prennent davantage de temps à la main, sans calculatrice. Il faut s'attendre à laisser quelques cases vides. La feuille récompense l'ordre dans lequel on traite les questions autant que la vitesse~; on commence donc par cueillir celles dont la réponse saute aux yeux.

\textbf{Où l'on perd des points.} Repérer la famille de chaque question manquée ou passée.
% Generated by code/interviews-1/03-speed-drills/python/gen_drill.py. Do not edit. French copy by draft/pilot-quant-fr/ivgen_fr.py.
\ivfamily{Additionner et soustraire de grands nombres}{2, 6, 11, 15, 17, 28, 30, 47}{Arrondir un nombre puis corriger~: $8142 - 3678 = 8142 - 3700 + 22 = 4464$.}
\ivfamily{Deux chiffres fois deux chiffres}{13, 16, 18, 19, 24, 29, 37, 39, 45, 50}{Décomposer un facteur, comme dans $37 \times 46 = 37 \times 40 + 37 \times 6 = 1480 + 222 = 1702$. Si les facteurs sont proches, utiliser plutôt la différence de deux carrés.}
\ivfamily{Carrés}{12, 20, 21, 36, 40, 48}{$n^2 = (n - d)(n + d) + d^2$, avec $d$ choisi pour rendre un facteur rond~: $48^2 = 50 \times 46 + 4 = 2304$. Un nombre finissant par 5~: $65^2$ vaut $6 \times 7$ suivi de 25.}
\ivfamily{Autour d'un nombre rond}{9, 31, 32, 35, 42}{$103 \times 97 = 100^2 - 3^2$, et $998 \times 6 = 6000 - 12$.}
\ivfamily{Pourcentages de n'importe quel nombre}{3, 4, 7, 25, 27, 34, 43}{Prendre 1~\% ou 10~\% et monter, ou permuter les deux~: $24\%$ de 350 égale $350\%$ de 24, soit 84.}
\ivfamily{Fractions}{1, 14, 22, 23, 33, 38, 49}{Utiliser le plus petit dénominateur commun. Pour diviser, multiplier par la fraction inverse, et simplifier d'abord.}
\ivfamily{Réponses décimales}{5, 8, 10, 26, 41, 44, 46}{Connaître les huitièmes et les seizièmes~: $\tfrac18 = 0.125$ et $\tfrac1{16} = 0.0625$. Pour diviser par 40, diviser par 4 puis par 10.}


\iqwrong{Ralentir pour être sûr de chaque question, comme sur la feuille 1. Ici, une erreur ne coûte qu'un point. Une feuille avec 30 bonnes et 2 mauvaises donne 28, et 40 bonnes avec 12 mauvaises donnent aussi 28~: à ce taux d'erreur, aller plus vite et moins soigneusement ne fait donc rien perdre. Sur cette feuille, la vitesse compte davantage que sur la première.}
\iqfollow{Puis-je imprimer la clé avec une autre répartition~?}{Oui. On lance \texttt{gen\_drill.py} avec le numéro de feuille 2 et une autre graine. Il écrit de nouvelles questions et une nouvelle clé, tirées des mêmes familles.}
\end{solution}

\begin{solution}{iq:iv1:speed-drills:3}
\textbf{Feuille 3.} Corriger la feuille avec cette clé, qui donne la lettre et la valeur. Le score est deux fois le nombre de bonnes réponses moins le nombre de mauvaises.
\par\nopagebreak
% Generated by code/interviews-1/03-speed-drills/python/gen_drill.py. Do not edit. French copy by draft/pilot-quant-fr/ivgen_fr.py.
\begin{ivdrillkey}
\ivkey{1}{$\mathrm{D}\;68$} & \ivkey{9}{$\mathrm{D}\;92$} & \ivkey{17}{$\mathrm{A}\;9919$} & \ivkey{25}{$\mathrm{B}\;80$} & \ivkey{33}{$\mathrm{A}\;57393$} \\
\ivkey{2}{$\mathrm{C}\;21522$} & \ivkey{10}{$\mathrm{B}\;140.3$} & \ivkey{18}{$\mathrm{A}\;72$} & \ivkey{26}{$\mathrm{C}\;9964$} & \ivkey{34}{$\mathrm{D}\;385$} \\
\ivkey{3}{$\mathrm{C}\;9.12$} & \ivkey{11}{$\mathrm{D}\;30.87$} & \ivkey{19}{$\mathrm{D}\;6889$} & \ivkey{27}{$\mathrm{A}\;20010$} & \ivkey{35}{$\mathrm{C}\;18.2$} \\
\ivkey{4}{$\mathrm{A}\;48.75$} & \ivkey{12}{$\mathrm{B}\;52.9$} & \ivkey{20}{$\mathrm{B}\;26790$} & \ivkey{28}{$\mathrm{C}\;180.9$} & \ivkey{36}{$\mathrm{A}\;97$} \\
\ivkey{5}{$\mathrm{B}\;455$} & \ivkey{13}{$\mathrm{C}\;8228$} & \ivkey{21}{$\mathrm{B}\;8464$} & \ivkey{29}{$\mathrm{D}\;31.36$} &  \\
\ivkey{6}{$\mathrm{B}\;50.7$} & \ivkey{14}{$\mathrm{C}\;12160$} & \ivkey{22}{$\mathrm{C}\;60$} & \ivkey{30}{$\mathrm{C}\;34.2$} &  \\
\ivkey{7}{$\mathrm{D}\;67$} & \ivkey{15}{$\mathrm{D}\;43169$} & \ivkey{23}{$\mathrm{A}\;82.8$} & \ivkey{31}{$\mathrm{D}\;35$} &  \\
\ivkey{8}{$\mathrm{C}\;26.88$} & \ivkey{16}{$\mathrm{D}\;2916$} & \ivkey{24}{$\mathrm{A}\;44400$} & \ivkey{32}{$\mathrm{A}\;210$} &  \\
\end{ivdrillkey}


\textbf{Quand laisser une case vide.} Avec $+2$ et $-1$, une réponse devinée, juste avec la probabilité $p$, ajoute en moyenne $2p - (1 - p) = 3p - 1$ au score, de sorte qu'elle paie lorsque $p > \tfrac13$. Un tirage au hasard parmi quatre options fait perdre $\tfrac14$ de point. Parmi trois, il ne rapporte rien. Parmi deux, il rapporte $\tfrac12$. On devine donc dès que deux options ont été éliminées, et l'on ne s'arrête pour calculer que si ce n'est pas le cas.

\textbf{Le rythme.} Six minutes pour 36 questions font 10 secondes par question, ce qui est trop peu pour la plupart des calculs complets. On travaille plutôt sur les options. Chaque question comporte trois mauvaises options~: une proche, avec le mauvais chiffre des unités~; une avec le bon chiffre des unités et une erreur de dizaines~; et le plus souvent une très éloignée de la bonne taille. Un contrôle du chiffre des unités élimine la première, l'arrondi élimine la plus éloignée, ce qui laisse un choix entre deux pour quelques secondes. La somme des chiffres, prise modulo 9, tranche alors souvent.

\textbf{Où l'on perd des points.} Repérer la famille de chaque question manquée ou passée.
% Generated by code/interviews-1/03-speed-drills/python/gen_drill.py. Do not edit. French copy by draft/pilot-quant-fr/ivgen_fr.py.
\ivfamily{Trois chiffres fois deux chiffres}{2, 13, 14, 15, 20, 24, 27, 33}{Ne pas finir le calcul. Vérifier le chiffre des unités, puis l'ordre de grandeur, puis la somme des chiffres. Pour $378 \times 46$, le chiffre des unités est celui de $8 \times 6$, et l'ordre de grandeur est voisin de $400 \times 45$.}
\ivfamily{Carrés et quasi-carrés}{16, 17, 19, 21, 26}{$67^2$ finit par 9, et est proche de $70 \times 64$. $104 \times 96$ vaut $10000 - 16$.}
\ivfamily{Division exacte}{1, 7, 18, 25, 31, 36}{Vérifier en calcul mental~: le quotient multiplié par le diviseur doit finir par le bon chiffre. Son ordre de grandeur vient de l'arrondi du diviseur.}
\ivfamily{Pourcentages}{6, 10, 12, 23, 28, 30}{$17\%$ de 640 vaut $10\% + 5\% + 2\%$, soit $64 + 32 + 12.8$.}
\ivfamily{Produits de décimaux}{3, 4, 8, 11, 29, 35}{Compter d'abord les décimales. Les propositions ne diffèrent alors que par des chiffres qu'on peut tester~: $3.4 \times 1.7$ a deux décimales et finit par 8.}
\ivfamily{Fractions d'un nombre}{5, 9, 22, 32, 34}{Diviser d'abord par le dénominateur, puis multiplier par le numérateur. Pour $\tfrac78$, retrancher un huitième.}


\iqwrong{Calculer chaque question en entier et ne vérifier les options qu'à la fin. Vingt questions traitées en entier et seize laissées vides donnent 40. En travaillant sur les options, 24 bonnes et 12 mauvaises donnent 36~: à ce taux d'erreur, les deux méthodes sont proches, mais une personne qui élimine bien les options obtient bien plus de 24 bonnes réponses.}
\iqfollow{Et si le vrai test ne propose pas d'options~?}{Alors les contrôles de la feuille 3 deviennent le moyen de repérer les erreurs~: le chiffre des unités, la taille et la somme des chiffres s'appliquent toujours. On lance \texttt{gen\_drill.py} avec le numéro de feuille 3 et une autre graine pour obtenir une nouvelle feuille.}
\end{solution}

\begin{solution}{iq:iv1:speed-drills:4}
\textbf{Feuille 4.} Corriger la feuille avec cette clé. Le score est le nombre de bonnes réponses moins le nombre de mauvaises.
\par\nopagebreak
% Generated by code/interviews-1/03-speed-drills/python/gen_drill.py. Do not edit. French copy by draft/pilot-quant-fr/ivgen_fr.py.
\begin{ivdrillkey}
\ivkey{1}{$240$} & \ivkey{12}{$\frac{7}{144}$}\rule[-9.5pt]{0pt}{23pt} & \ivkey{23}{$65$} & \ivkey{34}{$765$} & \ivkey{45}{$\frac{23}{125}$}\rule[-9.5pt]{0pt}{23pt} \\
\ivkey{2}{$77.5$} & \ivkey{13}{$2.5$} & \ivkey{24}{$92$} & \ivkey{35}{$40$} & \ivkey{46}{$63$} \\
\ivkey{3}{$92.5$} & \ivkey{14}{$\frac{9}{20}$}\rule[-9.5pt]{0pt}{23pt} & \ivkey{25}{$40$} & \ivkey{36}{$\frac{39}{40}$}\rule[-9.5pt]{0pt}{23pt} & \ivkey{47}{$78$} \\
\ivkey{4}{$612$} & \ivkey{15}{$85$} & \ivkey{26}{$\frac{7}{20}$}\rule[-9.5pt]{0pt}{23pt} & \ivkey{37}{$\frac{1}{25}$}\rule[-9.5pt]{0pt}{23pt} & \ivkey{48}{$\frac{1}{2}$}\rule[-9.5pt]{0pt}{23pt} \\
\ivkey{5}{$60$} & \ivkey{16}{$365.5$} & \ivkey{27}{$\frac{13}{125}$}\rule[-9.5pt]{0pt}{23pt} & \ivkey{38}{$8.4$} & \ivkey{49}{$91.2$} \\
\ivkey{6}{$665$} & \ivkey{17}{$178.5$} & \ivkey{28}{$\frac{1}{125}$}\rule[-9.5pt]{0pt}{23pt} & \ivkey{39}{$14$} & \ivkey{50}{$\frac{19}{24}$}\rule[-9.5pt]{0pt}{23pt} \\
\ivkey{7}{$\frac{6}{7}$}\rule[-9.5pt]{0pt}{23pt} & \ivkey{18}{$214.5$} & \ivkey{29}{$260$} & \ivkey{40}{$55$} & \ivkey{51}{$1050$} \\
\ivkey{8}{$552$} & \ivkey{19}{$15$} & \ivkey{30}{$\frac{29}{40}$}\rule[-9.5pt]{0pt}{23pt} & \ivkey{41}{$8$} & \ivkey{52}{$10$} \\
\ivkey{9}{$\frac{1}{6}$}\rule[-9.5pt]{0pt}{23pt} & \ivkey{20}{$95$} & \ivkey{31}{$99$} & \ivkey{42}{$\frac{13}{35}$}\rule[-9.5pt]{0pt}{23pt} & \ivkey{53}{$156$} \\
\ivkey{10}{$\frac{26}{21}$}\rule[-9.5pt]{0pt}{23pt} & \ivkey{21}{$5$} & \ivkey{32}{$360$} & \ivkey{43}{$714$} & \ivkey{54}{$\frac{9}{7}$}\rule[-9.5pt]{0pt}{23pt} \\
\ivkey{11}{$35$} & \ivkey{22}{$50$} & \ivkey{33}{$\frac{71}{125}$}\rule[-9.5pt]{0pt}{23pt} & \ivkey{44}{$\frac{1}{20}$}\rule[-9.5pt]{0pt}{23pt} &  \\
\end{ivdrillkey}


\textbf{Quand laisser une case vide.} Avec $+1$ et $-1$, une réponse devinée, juste avec la probabilité $p$, ajoute en moyenne $2p - 1$ au score, de sorte qu'elle paie lorsque $p > \tfrac12$.

\textbf{Le rythme.} Six minutes pour 54 questions font moins de 7 secondes par question. Cette feuille porte sur le passage d'une écriture à une autre d'un même nombre, donc les questions les plus rapides sont celles que l'on connaît par cœur. On les traite en premier, puis l'on revient aux autres.

\textbf{Où l'on perd des points.} Repérer la famille de chaque question manquée ou passée.
% Generated by code/interviews-1/03-speed-drills/python/gen_drill.py. Do not edit. French copy by draft/pilot-quant-fr/ivgen_fr.py.
\ivfamily{Fractions en pourcentages}{2, 3, 13, 20, 21, 24, 39, 40, 41, 47}{Connaître les huitièmes et les seizièmes. Pour un dénominateur de 20, 25 ou 50, ramener à 100~: $\tfrac{7}{25} = \tfrac{28}{100}$. Pour 40, prendre la moitié de la valeur sur 20.}
\ivfamily{Décimaux en fractions}{26, 27, 28, 30, 33, 36, 37, 45}{Lire le décimal sur une puissance de 10, puis simplifier~: $0.375$ vaut $\tfrac{375}{1000}$, soit $\tfrac38$. La plupart sont des multiples de $\tfrac18$, $\tfrac1{20}$, $\tfrac1{25}$ ou $\tfrac1{40}$.}
\ivfamily{Variation en pourcentage}{1, 4, 6, 8, 29, 32, 34, 43, 51, 53}{Ajouter ou soustraire le pourcentage de la base, ou multiplier directement~: $15\%$ en hausse s'écrit $\times 1.15$, et $20\%$ en baisse s'écrit $\times 0.8$.}
\ivfamily{En pourcentage de}{11, 15, 19, 22, 23, 25, 35, 52}{Simplifier d'abord, puis ramener à 100~: $36$ sur $120$ vaut $\tfrac{3}{10}$, soit $30\%$.}
\ivfamily{Pourcentages de n'importe quel nombre}{5, 16, 17, 18, 31, 38, 46, 49}{Prendre 1~\% ou 10~\% et monter, ou permuter les deux~: $p\%$ de $n$ égale $n\%$ de $p$.}
\ivfamily{Fractions}{7, 9, 10, 12, 14, 42, 44, 48, 50, 54}{Utiliser le plus petit dénominateur commun. Pour diviser, multiplier par la fraction inverse, et simplifier d'abord.}


\iqwrong{Faire une division longue pour chaque fraction. Les huitièmes, les vingtièmes et les quarantièmes reviennent partout sur cette feuille, et chacun que l'on connaît fait gagner quelques secondes sur toutes les feuilles suivantes.}
\iqfollow{Que faut-il apprendre par cœur en premier~?}{Les huitièmes~: $12.5$, $37.5$, $62.5$ et $87.5$. Ensuite les seizièmes, qui en sont la moitié~: $\tfrac{3}{16}$ vaut $18.75~\%$.}
\end{solution}

\begin{solution}{iq:iv1:speed-drills:5}
\textbf{Feuille 5.} Corriger la feuille avec cette clé. Le score est le nombre de bonnes réponses.
\par\nopagebreak
% Generated by code/interviews-1/03-speed-drills/python/gen_drill.py. Do not edit. French copy by draft/pilot-quant-fr/ivgen_fr.py.
\begin{ivdrillkey}
\ivkey{1}{$90$} & \ivkey{19}{$6$} & \ivkey{37}{$84$} & \ivkey{55}{$1736$} & \ivkey{73}{$41$} \\
\ivkey{2}{$60$} & \ivkey{20}{$91$} & \ivkey{38}{$382$} & \ivkey{56}{$128$} & \ivkey{74}{$91$} \\
\ivkey{3}{$20$} & \ivkey{21}{$210$} & \ivkey{39}{$92$} & \ivkey{57}{$1050$} & \ivkey{75}{$42$} \\
\ivkey{4}{$52$} & \ivkey{22}{$20$} & \ivkey{40}{$487$} & \ivkey{58}{$136$} & \ivkey{76}{$40$} \\
\ivkey{5}{$48$} & \ivkey{23}{$144$} & \ivkey{41}{$831$} & \ivkey{59}{$171$} & \ivkey{77}{$460$} \\
\ivkey{6}{$119$} & \ivkey{24}{$876$} & \ivkey{42}{$456$} & \ivkey{60}{$366$} & \ivkey{78}{$68$} \\
\ivkey{7}{$153$} & \ivkey{25}{$520$} & \ivkey{43}{$481$} & \ivkey{61}{$64$} & \ivkey{79}{$144$} \\
\ivkey{8}{$8$} & \ivkey{26}{$112$} & \ivkey{44}{$363$} & \ivkey{62}{$356$} & \ivkey{80}{$84$} \\
\ivkey{9}{$1134$} & \ivkey{27}{$68$} & \ivkey{45}{$61$} & \ivkey{63}{$85$} & \ivkey{81}{$190$} \\
\ivkey{10}{$844$} & \ivkey{28}{$1003$} & \ivkey{46}{$104$} & \ivkey{64}{$650$} & \ivkey{82}{$98$} \\
\ivkey{11}{$672$} & \ivkey{29}{$56$} & \ivkey{47}{$90$} & \ivkey{65}{$942$} & \ivkey{83}{$750$} \\
\ivkey{12}{$552$} & \ivkey{30}{$29$} & \ivkey{48}{$72$} & \ivkey{66}{$95$} & \ivkey{84}{$1128$} \\
\ivkey{13}{$74$} & \ivkey{31}{$36$} & \ivkey{49}{$539$} & \ivkey{67}{$402$} & \ivkey{85}{$726$} \\
\ivkey{14}{$1368$} & \ivkey{32}{$42$} & \ivkey{50}{$827$} & \ivkey{68}{$870$} & \ivkey{86}{$504$} \\
\ivkey{15}{$333$} & \ivkey{33}{$108$} & \ivkey{51}{$35$} & \ivkey{69}{$476$} & \ivkey{87}{$532$} \\
\ivkey{16}{$450$} & \ivkey{34}{$215$} & \ivkey{52}{$478$} & \ivkey{70}{$656$} & \ivkey{88}{$64$} \\
\ivkey{17}{$847$} & \ivkey{35}{$57$} & \ivkey{53}{$427$} & \ivkey{71}{$433$} & \ivkey{89}{$370$} \\
\ivkey{18}{$80$} & \ivkey{36}{$560$} & \ivkey{54}{$152$} & \ivkey{72}{$399$} & \ivkey{90}{$348$} \\
\end{ivdrillkey}


\textbf{Quand laisser une case vide.} Sans pénalité, une réponse devinée ne coûte rien~: il ne faut donc jamais laisser de case vide. Si le temps est sur le point de s'écouler, on écrit la valeur la plus proche dont on dispose.

\textbf{Le rythme.} Sept minutes pour 90 questions font moins de 5 secondes par question. C'est trop rapide pour que la plupart des gens finissent, quelle que soit la simplicité des questions. La feuille teste si l'arithmétique simple est automatique, car la moindre réflexion, même brève, consomme le temps.

\textbf{Où l'on perd des points.} Repérer la famille de chaque question manquée ou passée.
% Generated by code/interviews-1/03-speed-drills/python/gen_drill.py. Do not edit. French copy by draft/pilot-quant-fr/ivgen_fr.py.
\ivfamily{Additionner et soustraire}{10, 11, 17, 24, 25, 28, 34, 36, 38, 40, 41, 42, 44, 49, 50, 53, 64, 66, 68, 83}{Travailler depuis la gauche, et ne pas écrire les retenues.}
\ivfamily{Tables de multiplication jusqu'à 19}{1, 2, 4, 5, 6, 7, 18, 20, 23, 26, 27, 31, 32, 33, 35, 37, 46, 48, 56, 58, 59, 63, 79, 80, 82}{Décomposer les nombres de 11 à 19~: $17 \times 8 = 80 + 56$.}
\ivfamily{Multiplier par un chiffre}{9, 12, 14, 15, 16, 54, 55, 57, 65, 70, 72, 84, 85, 86, 89}{Décomposer le plus grand nombre, les centaines d'abord.}
\ivfamily{Doubler et diviser par deux}{43, 52, 60, 62, 67, 69, 71, 77, 87, 90}{Procéder chiffre par chiffre depuis la gauche. Cela ne demande aucune réflexion.}
\ivfamily{Division exacte}{13, 30, 39, 45, 51, 73, 74, 75, 78, 88}{Penser la table de multiplication à l'envers, et décomposer un long nombre en parties faciles~: $476 \div 7 = 420 \div 7 + 56 \div 7$.}
\ivfamily{Pourcentages faciles}{3, 8, 19, 21, 22, 29, 47, 61, 76, 81}{$10\%$ déplace la virgule, $20\%$ double ce résultat, $25\%$ est un quart et $50\%$ une moitié.}


\iqwrong{Relire son travail. Sans pénalité, une relecture ne rapporte rien et coûte du temps. Une séance de 70 questions toutes justes donne 70, et une de 90 questions avec 12 mauvaises donne 78. Lorsqu'une mauvaise réponse coûte $-1$ ou plus pour chaque $+1$, les calculs changent (avec $+1$ et $-2$, 70 contre 54)~: il faut donc lire le barème avant de commencer.}
\iqfollow{Et si je finis avec du temps devant moi~?}{Alors ce niveau est automatisé, ce qui est bon à savoir. Il faut reprendre les questions passées, puis passer aux feuilles plus difficiles.}
\end{solution}

\enlargethispage{-3\baselineskip}%
\begin{solution}{iq:iv1:speed-drills:6}
\textbf{Feuille 6.} Corriger la feuille avec cette clé. Le score est trois fois le nombre de bonnes réponses moins le nombre de mauvaises.
\par\nopagebreak
% Generated by code/interviews-1/03-speed-drills/python/gen_drill.py. Do not edit. French copy by draft/pilot-quant-fr/ivgen_fr.py.
\begin{ivdrillkey}
\ivkey{1}{$\mathrm{D}\;9.48$} & \ivkey{9}{$\mathrm{B}\;9936$} & \ivkey{17}{$\mathrm{D}\;35$} & \ivkey{25}{$\mathrm{C}\;20.79$} & \ivkey{33}{$\mathrm{C}\;140$} \\
\ivkey{2}{$\mathrm{B}\;240.9$} & \ivkey{10}{$\mathrm{B}\;7938$} & \ivkey{18}{$\mathrm{D}\;62.1$} & \ivkey{26}{$\mathrm{B}\;462$} & \ivkey{34}{$\mathrm{D}\;9951$} \\
\ivkey{3}{$\mathrm{C}\;69126$} & \ivkey{11}{$\mathrm{C}\;66$} & \ivkey{19}{$\mathrm{A}\;58.52$} & \ivkey{27}{$\mathrm{A}\;81.4$} & \ivkey{35}{$\mathrm{D}\;35$} \\
\ivkey{4}{$\mathrm{D}\;21.32$} & \ivkey{12}{$\mathrm{D}\;455$} & \ivkey{20}{$\mathrm{A}\;476$} & \ivkey{28}{$\mathrm{C}\;44.24$} & \ivkey{36}{$\mathrm{B}\;56$} \\
\ivkey{5}{$\mathrm{A}\;13320$} & \ivkey{13}{$\mathrm{A}\;56$} & \ivkey{21}{$\mathrm{D}\;33927$} & \ivkey{29}{$\mathrm{B}\;66$} &  \\
\ivkey{6}{$\mathrm{D}\;49771$} & \ivkey{14}{$\mathrm{B}\;36.1$} & \ivkey{22}{$\mathrm{D}\;7225$} & \ivkey{30}{$\mathrm{C}\;61$} &  \\
\ivkey{7}{$\mathrm{C}\;37$} & \ivkey{15}{$\mathrm{D}\;14534$} & \ivkey{23}{$\mathrm{B}\;170.2$} & \ivkey{31}{$\mathrm{D}\;30120$} &  \\
\ivkey{8}{$\mathrm{D}\;7.14$} & \ivkey{16}{$\mathrm{A}\;72$} & \ivkey{24}{$\mathrm{D}\;6724$} & \ivkey{32}{$\mathrm{D}\;3844$} &  \\
\end{ivdrillkey}


\textbf{Quand laisser une case vide.} Avec $+3$ et $-1$, une réponse devinée, juste avec la probabilité $p$, ajoute en moyenne $4p - 1$ au score, de sorte qu'elle paie lorsque $p > \tfrac14$. Un tirage au hasard parmi quatre options rapporte exactement 0. Si l'on élimine une option, un tirage parmi trois rapporte $\tfrac13$. Si l'on en élimine deux, un tirage entre deux rapporte 1. Il ne faut donc jamais laisser une question vide dès que l'on peut éliminer ne serait-ce qu'une option.

\textbf{Le rythme.} Cinq minutes pour 36 questions font environ 8 secondes par question. Comme sur la feuille 3, chaque question comporte une mauvaise option proche, avec le mauvais chiffre des unités, une avec une erreur de dizaines, et le plus souvent une de mauvaise taille.

\textbf{Où l'on perd des points.} Repérer la famille de chaque question manquée ou passée.
% Generated by code/interviews-1/03-speed-drills/python/gen_drill.py. Do not edit. French copy by draft/pilot-quant-fr/ivgen_fr.py.
\ivfamily{Trois chiffres fois deux chiffres}{3, 5, 6, 10, 15, 21, 31}{Tester d'abord le chiffre des unités, puis l'ordre de grandeur.}
\ivfamily{Carrés et quasi-carrés}{9, 22, 24, 32, 34}{Vérifier le chiffre des unités, puis arrondir. $(100 + d)(100 - d)$ vaut $10000 - d^2$.}
\ivfamily{Division longue}{7, 13, 16, 17, 29, 30, 35}{Vérifier en multipliant~: le chiffre des unités du quotient multiplié par celui du diviseur doit donner le chiffre des unités du nombre.}
\ivfamily{Pourcentages}{2, 14, 18, 23, 27, 36}{Construire le pourcentage à partir de 10~\%, 5~\% et 1~\%, et s'arrêter dès qu'il ne reste qu'une seule option.}
\ivfamily{Produits de décimaux}{1, 4, 8, 19, 25, 28}{Compter les décimales, puis tester le chiffre des unités du produit.}
\ivfamily{Fractions d'un nombre}{11, 12, 20, 26, 33}{Diviser d'abord par le dénominateur, puis multiplier par le numérateur.}


\iqwrong{Laisser des cases vides par prudence, comme si la pénalité était lourde. Éliminer une option sur chacune des 36 questions puis deviner rapporte 12 en moyenne, contre 0 pour les cases vides. La règle ici récompense toute information, même petite.}
\iqfollow{Pourquoi le tirage au hasard vaut-il exactement zéro~?}{Parce que la règle a été conçue ainsi~: avec quatre options, $+3$ et $-1$ se compensent exactement. Si un test fixe sa pénalité de cette façon, on ne perd rien à deviner, et l'on gagne grâce à tout ce que l'on sait.}
\end{solution}

\begin{solution}{iq:iv1:speed-drills:7}
\textbf{Feuille 7.} Corriger la feuille avec cette clé. Une réponse est juste si elle se trouve à moins de 5~\% de la valeur de la clé. La clé donne la valeur exacte à trois chiffres significatifs, ce qui suffit ici. Le score est le nombre de bonnes réponses moins le nombre de mauvaises.
\par\nopagebreak
% Generated by code/interviews-1/03-speed-drills/python/gen_drill.py. Do not edit. French copy by draft/pilot-quant-fr/ivgen_fr.py.
\begin{ivdrillkey}
\ivkey{1}{$85.3$} & \ivkey{7}{$1.34$} & \ivkey{13}{$1.65$} & \ivkey{19}{$76.2$} & \ivkey{25}{$465$} \\
\ivkey{2}{$13.2$} & \ivkey{8}{$97.9$} & \ivkey{14}{$9.84$} & \ivkey{20}{$17.6$} & \ivkey{26}{$1720$} \\
\ivkey{3}{$26.5$} & \ivkey{9}{$4050$} & \ivkey{15}{$1.28$} & \ivkey{21}{$1420$} & \ivkey{27}{$3.91$} \\
\ivkey{4}{$909$} & \ivkey{10}{$855$} & \ivkey{16}{$2.65$} & \ivkey{22}{$8.95$} & \ivkey{28}{$3440$} \\
\ivkey{5}{$6.68$} & \ivkey{11}{$17.6$} & \ivkey{17}{$2560$} & \ivkey{23}{$6.21$} & \ivkey{29}{$60.2$} \\
\ivkey{6}{$1070$} & \ivkey{12}{$1.63$} & \ivkey{18}{$7120$} & \ivkey{24}{$1.11$} & \ivkey{30}{$1.47$} \\
\end{ivdrillkey}


\textbf{Quand laisser une case vide.} Avec $+1$ et $-1$, une réponse devinée, juste avec la probabilité $p$, ajoute en moyenne $2p - 1$ au score, de sorte qu'elle paie lorsque $p > \tfrac12$. Une estimation grossière, dont on est sûr à 5~\% près, passe largement ce seuil. Si l'on n'a que l'ordre de grandeur, on laisse la case vide.

\textbf{Le rythme.} Cinq minutes pour 30 questions font 10 secondes par question, et un calcul exact prendrait presque tout ce temps pour un seul produit. On arrondit à des nombres que l'on sait multiplier, et l'on ne corrige que si les arrondis vont dans le même sens.

\textbf{Où l'on perd des points.} Repérer la famille de chaque question manquée ou passée.
% Generated by code/interviews-1/03-speed-drills/python/gen_drill.py. Do not edit. French copy by draft/pilot-quant-fr/ivgen_fr.py.
\ivfamily{Produits}{6, 10, 17, 18, 21, 28}{Arrondir un facteur vers le haut et l'autre vers le bas, de la même proportion environ~: $38.7 \times 41.2$ est proche de $40 \times 40$.}
\ivfamily{Quotients}{2, 3, 5, 14, 20, 22}{Arrondir le diviseur à un nombre facile, puis corriger du pourcentage dont on l'a déplacé.}
\ivfamily{Racines carrées}{1, 8, 11, 19, 29}{Écrire le nombre comme un carré multiplié par un petit facteur, et déplacer la racine de la moitié de l'écart en pourcentage~: $\sqrt{5300}$ vaut environ $70 + \tfrac{400}{140}$.}
\ivfamily{Intérêts composés}{7, 12, 15, 16, 30}{Prendre $1.05^{10} \approx e^{0.5}$, et retrancher un peu, car le logarithme de $1 + r$ est inférieur à $r$.}
\ivfamily{Exponentielles et logarithmes}{13, 23, 24, 27}{Utiliser $e^x \approx 1 + x + \tfrac{x^2}{2}$ pour un petit $x$. Construire chaque logarithme à partir de $\ln 2$ et $\ln 10$.}
\ivfamily{Pourcentages}{4, 9, 25, 26}{Arrondir le nombre, et non le pourcentage, puis corriger du même petit facteur.}


\iqwrong{Calcul exact. Le produit $38.7 \times 41.2$ vaut $1594.44$, et $40 \times 40 = 1600$ se calcule en une seconde et ne s'en écarte que de $0.35~\%$. Arrondir les deux facteurs vers le haut, en revanche, additionne les deux erreurs~: $40 \times 45$ pour $38.7 \times 42.6$ est trop grand de $9~\%$.}
\iqfollow{Que faut-il arrondir en premier~?}{Arrondissez dans des sens opposés si vous le pouvez, car les erreurs se compensent. Sinon, notez l'ampleur de l'erreur et retranchez ce pourcentage au résultat.}
\end{solution}

\enlargethispage{-1\baselineskip}%
\begin{solution}{iq:iv1:speed-drills:8}
\textbf{Feuille 8.} Corriger la feuille avec cette clé. Votre score est le double du nombre de bonnes réponses moins le nombre de mauvaises.
\par\nopagebreak
% Generated by code/interviews-1/03-speed-drills/python/gen_drill.py. Do not edit. French copy by draft/pilot-quant-fr/ivgen_fr.py.
\begin{ivdrillkey}
\ivkey{1}{$25$} & \ivkey{9}{$17$} & \ivkey{17}{$13$} & \ivkey{25}{$8$} & \ivkey{33}{$19$} \\
\ivkey{2}{$17$} & \ivkey{10}{$6$} & \ivkey{18}{$9$} & \ivkey{26}{$4$} & \ivkey{34}{$3$} \\
\ivkey{3}{$14$} & \ivkey{11}{$10$} & \ivkey{19}{$5$} & \ivkey{27}{$16$} & \ivkey{35}{$90$} \\
\ivkey{4}{$14$} & \ivkey{12}{$7$} & \ivkey{20}{$11$} & \ivkey{28}{$25$} & \ivkey{36}{$6$} \\
\ivkey{5}{$105$} & \ivkey{13}{$88$} & \ivkey{21}{$9$} & \ivkey{29}{$18$} & \ivkey{37}{$14$} \\
\ivkey{6}{$4$} & \ivkey{14}{$10$} & \ivkey{22}{$5$} & \ivkey{30}{$13$} & \ivkey{38}{$8$} \\
\ivkey{7}{$18$} & \ivkey{15}{$21$} & \ivkey{23}{$14$} & \ivkey{31}{$22$} & \ivkey{39}{$45$} \\
\ivkey{8}{$8$} & \ivkey{16}{$5$} & \ivkey{24}{$12$} & \ivkey{32}{$36$} & \ivkey{40}{$75$} \\
\end{ivdrillkey}


\textbf{Quand laisser une case vide.} Avec $+2$ et $-1$, une réponse devinée, juste avec la probabilité $p$, ajoute en moyenne $3p - 1$, de sorte qu'elle paie lorsque $p > \tfrac13$. Chaque réponse ici est un nombre entier, donc une fraction signale une erreur d'inattention. Refaire le calcul s'il est rapide~; sinon, laisser la case vide, car l'écrire revient à coup sûr à $-1$.

\textbf{Le rythme.} Cinq minutes pour 40 questions font $7.5$ secondes par question. N'écrivez pas les étapes. Défaites les opérations de tête, et ne remettez la valeur dans l'équation que si vous en avez le temps.

\textbf{Où l'on perd des points.} Repérer la famille de chaque question manquée ou passée.
% Generated by code/interviews-1/03-speed-drills/python/gen_drill.py. Do not edit. French copy by draft/pilot-quant-fr/ivgen_fr.py.
\ivfamily{Deux étapes}{7, 9, 15, 23, 24, 26, 27, 28, 31, 38}{Défaire l'addition, puis la multiplication, mentalement.}
\ivfamily{Parenthèses}{8, 11, 19, 21, 22, 30, 33, 37}{Diviser d'abord par le facteur, s'il divise exactement, et laisser la parenthèse telle quelle.}
\ivfamily{Deux membres}{3, 4, 6, 12, 14, 16, 17, 36}{Regrouper les termes en $x$ d'un côté, et les nombres de l'autre.}
\ivfamily{Avec une fraction}{5, 10, 13, 25, 32, 35, 39, 40}{Éliminer la fraction en multipliant, ou lire $\tfrac{12}{x} = 1.5$ comme $x = \tfrac{12}{1.5}$.}
\ivfamily{Puissances}{1, 2, 18, 20, 29, 34}{Connaître les carrés jusqu'à 25 et les puissances de 2 et de 3.}


\iqwrong{Développer chaque parenthèse. Pour $5(x - 2) = 35$, diviser d'abord par 5 et obtenir $x - 2 = 7$. Développer donne $5x - 10 = 35$, soit une étape de plus pour le même résultat.}
\iqfollow{Une vérification par substitution vaut-elle le temps qu'elle prend~?}{Seulement si elle est rapide. Pour $7x - 9 = 3x + 15$, la valeur $x = 6$ donne 33 des deux côtés, et cela prend à peu près autant de temps que de résoudre. Sur une feuille chronométrée, ne vérifier que les réponses qui semblent bizarres.}
\end{solution}

\begin{solution}{iq:iv1:speed-drills:9}
\textbf{Feuille 9.} Corriger la feuille avec cette clé. Votre score est le double du nombre de bonnes réponses moins trois fois le nombre de mauvaises.
\par\nopagebreak
% Generated by code/interviews-1/03-speed-drills/python/gen_drill.py. Do not edit. French copy by draft/pilot-quant-fr/ivgen_fr.py.
\begin{ivdrillkey}
\ivkey{1}{$120$} & \ivkey{10}{$4096$} & \ivkey{19}{$625$} & \ivkey{28}{$2401$} & \ivkey{37}{$5$} \\
\ivkey{2}{$79$} & \ivkey{11}{$324$} & \ivkey{20}{$5040$} & \ivkey{29}{$16384$} & \ivkey{38}{$512$} \\
\ivkey{3}{$69$} & \ivkey{12}{$2197$} & \ivkey{21}{$14641$} & \ivkey{30}{$8192$} & \ivkey{39}{$3375$} \\
\ivkey{4}{$289$} & \ivkey{13}{$42$} & \ivkey{22}{$441$} & \ivkey{31}{$13$} & \ivkey{40}{$90$} \\
\ivkey{5}{$10$} & \ivkey{14}{$4$} & \ivkey{23}{$2048$} & \ivkey{32}{$7$} & \ivkey{41}{$46$} \\
\ivkey{6}{$15625$} & \ivkey{15}{$4$} & \ivkey{24}{$66$} & \ivkey{33}{$24$} & \ivkey{42}{$43$} \\
\ivkey{7}{$6$} & \ivkey{16}{$72$} & \ivkey{25}{$6$} & \ivkey{34}{$2187$} & \ivkey{43}{$14$} \\
\ivkey{8}{$256$} & \ivkey{17}{$125$} & \ivkey{26}{$529$} & \ivkey{35}{$44$} & \ivkey{44}{$1331$} \\
\ivkey{9}{$361$} & \ivkey{18}{$362880$} & \ivkey{27}{$7$} & \ivkey{36}{$625$} & \ivkey{45}{$32768$} \\
\end{ivdrillkey}


\textbf{Quand laisser une case vide.} Avec $+2$ et $-3$, une réponse devinée, juste avec la probabilité $p$, ajoute en moyenne $5p - 3$, de sorte qu'elle ne paie que lorsque $p > \tfrac35$. La plupart des questions de cette feuille sont des faits que l'on connaît ou que l'on ignore. Laissez la case vide si vous n'êtes pas sûr.

\textbf{Le rythme.} Cinq minutes pour 45 questions font moins de 7 secondes par question. Cette feuille teste autant la mémoire que la vitesse. Tout ce qu'il faut calculer doit passer à un second tour.

\textbf{Où l'on perd des points.} Repérer la famille de chaque question manquée ou passée.
% Generated by code/interviews-1/03-speed-drills/python/gen_drill.py. Do not edit. French copy by draft/pilot-quant-fr/ivgen_fr.py.
\ivfamily{Puissances de deux}{10, 23, 29, 30, 38, 45}{$2^{10} = 1024$. À partir de là, doubler ou diviser par deux, ou lire $2^{13}$ comme $8 \times 1024$.}
\ivfamily{Autres puissances}{6, 12, 17, 19, 21, 28, 34, 39, 44}{Élever au carré, puis multiplier une fois de plus si besoin~: $7^4 = 49^2$. $11^3 = 1331$ et $5^n$ vaut $10^n$ divisé par $2^n$.}
\ivfamily{Carrés jusqu'à 29}{4, 8, 9, 11, 22, 26, 36}{Utiliser $n^2 = (n - d)(n + d) + d^2$, ou simplement les connaître.}
\ivfamily{Racines carrées}{2, 3, 13, 16, 24, 35, 41, 42}{La taille donne le chiffre des dizaines, et le chiffre des unités laisse au plus deux choix. Comparer avec le carré du nombre finissant par 5.}
\ivfamily{Racines cubiques}{14, 25, 31, 32, 43}{Connaître les cubes jusqu'à 15. Le chiffre des unités d'un cube fixe celui de sa racine.}
\ivfamily{Logarithmes}{5, 7, 15, 27, 37}{Réécrire le nombre comme une puissance de la base~: $\log_2 4096$ est la puissance de deux qui donne 4096.}
\ivfamily{Factorielles}{1, 18, 20, 33, 40}{Connaître jusqu'à $9! = 362\,880$. Un rapport tel que $\tfrac{10!}{8!}$ vaut simplement $10 \times 9$.}


\iqwrong{Construire $2^{13}$ en doublant à partir de 2. Partir de $2^{10} = 1024$, et $2^{13}$ vaut $8 \times 1024 = 8192$. Connaître un repère dans chaque suite suffit.}
\iqfollow{Lesquelles de ces valeurs valent-elles la peine d'être apprises par cœur~?}{Les puissances de 2 jusqu'à $2^{16} = 65\,536$, les carrés jusqu'à 29, les cubes jusqu'à 15 et les factorielles jusqu'à $9!$. Elles reviennent tout au long de la suite de ce livre.}
\end{solution}

\begin{solution}{iq:iv1:speed-drills:10}
\textbf{Feuille 10.} Corriger la feuille avec cette clé. Votre score est le nombre de bonnes réponses moins le nombre de mauvaises.
\par\nopagebreak
% Generated by code/interviews-1/03-speed-drills/python/gen_drill.py. Do not edit. French copy by draft/pilot-quant-fr/ivgen_fr.py.
\begin{ivdrillkey}
\ivkey{1}{$\frac{255}{256}$}\rule[-9.5pt]{0pt}{23pt} & \ivkey{9}{$120$} & \ivkey{17}{$55$} & \ivkey{25}{$\frac{1}{64}$}\rule[-9.5pt]{0pt}{23pt} & \ivkey{33}{$1287$} \\
\ivkey{2}{$90$} & \ivkey{10}{$56$} & \ivkey{18}{$62$} & \ivkey{26}{$66$} & \ivkey{34}{$\frac{5}{16}$}\rule[-9.5pt]{0pt}{23pt} \\
\ivkey{3}{$42$} & \ivkey{11}{$\frac{1}{512}$}\rule[-9.5pt]{0pt}{23pt} & \ivkey{19}{$33.5$} & \ivkey{27}{$\frac{63}{256}$}\rule[-9.5pt]{0pt}{23pt} & \ivkey{35}{$1$} \\
\ivkey{4}{$\frac{1023}{1024}$}\rule[-9.5pt]{0pt}{23pt} & \ivkey{12}{$23.5$} & \ivkey{20}{$336$} & \ivkey{28}{$18$} & \ivkey{36}{$84$} \\
\ivkey{5}{$\frac{2}{15}$}\rule[-9.5pt]{0pt}{23pt} & \ivkey{13}{$\frac{2}{7}$}\rule[-9.5pt]{0pt}{23pt} & \ivkey{21}{$30$} & \ivkey{29}{$39$} & \ivkey{37}{$\frac{15}{16}$}\rule[-9.5pt]{0pt}{23pt} \\
\ivkey{6}{$\frac{5}{32}$}\rule[-9.5pt]{0pt}{23pt} & \ivkey{14}{$78$} & \ivkey{22}{$182$} & \ivkey{30}{$\frac{1}{32}$}\rule[-9.5pt]{0pt}{23pt} & \ivkey{38}{$35$} \\
\ivkey{7}{$6$} & \ivkey{15}{$45$} & \ivkey{23}{$\frac{1}{32}$}\rule[-9.5pt]{0pt}{23pt} & \ivkey{31}{$\frac{21}{55}$}\rule[-9.5pt]{0pt}{23pt} & \ivkey{39}{$21$} \\
\ivkey{8}{$\frac{1}{26}$}\rule[-9.5pt]{0pt}{23pt} & \ivkey{16}{$\frac{6}{11}$}\rule[-9.5pt]{0pt}{23pt} & \ivkey{24}{$\frac{1}{21}$}\rule[-9.5pt]{0pt}{23pt} & \ivkey{32}{$\frac{5}{16}$}\rule[-9.5pt]{0pt}{23pt} & \ivkey{40}{$\frac{7}{64}$}\rule[-9.5pt]{0pt}{23pt} \\
\end{ivdrillkey}


\textbf{Quand laisser une case vide.} Avec $+1$ et $-1$, une réponse devinée ajoute en moyenne $2p - 1$, de sorte qu'elle paie lorsque $p > \tfrac12$.

\textbf{Le rythme.} Six minutes pour 40 questions font 9 secondes par question. C'est l'arithmétique qui sous-tend une réponse de probabilités. Ces questions sont du type où l'on se trompe autrement en fin de raisonnement juste.

\textbf{Où l'on perd des points.} Repérer la famille de chaque question manquée ou passée.
% Generated by code/interviews-1/03-speed-drills/python/gen_drill.py. Do not edit. French copy by draft/pilot-quant-fr/ivgen_fr.py.
\ivfamily{Coefficients binomiaux}{9, 10, 14, 15, 17, 26, 33, 36, 38, 39}{Simplifier avant de multiplier~: $\binom{12}{3} = \tfrac{12 \cdot 11 \cdot 10}{6} = 2 \cdot 11 \cdot 10$.}
\ivfamily{Rapports de factorielles}{2, 3, 20, 22}{Écrire seulement les facteurs restants~: $\tfrac{7!}{5!} = 7 \times 6$.}
\ivfamily{Puissances d'un demi}{1, 4, 11, 23, 25, 37}{$2^{-5} = \tfrac{1}{32}$. Le complément de $2^{-k}$ est $\tfrac{2^k - 1}{2^k}$.}
\ivfamily{Probabilités binomiales}{6, 27, 30, 32, 34, 40}{Simplifier d'abord les puissances de deux du coefficient, et donner le résultat sous forme irréductible.}
\ivfamily{Tirages successifs}{5, 8, 13, 16, 24, 31}{Simplifier en croix avant de multiplier~: $\tfrac59 \times \tfrac48 = \tfrac59 \times \tfrac12$.}
\ivfamily{Valeurs espérées}{7, 12, 18, 19, 21, 28, 29, 35}{Multiplier chaque probabilité par son montant, et placer la virgule en traitant $0.3$ comme 3 dixièmes.}


\iqwrong{Développer avant de simplifier. Pour $\binom{12}{3}$, calculer $12 \times 11 \times 10 = 1320$ puis diviser par 6 fonctionne, mais $2 \times 11 \times 10 = 220$ n'exige aucun grand nombre.}
\iqfollow{Comment vérifier rapidement un coefficient binomial~?}{Utiliser l'égalité $\binom{n}{k} = \binom{n - 1}{k - 1} \cdot \tfrac{n}{k}$. Par exemple, $\binom{10}{3} = \binom{9}{2} \cdot \tfrac{10}{3} = 36 \cdot \tfrac{10}{3} = 120$.}
\end{solution}

\begin{solution}{iq:iv1:speed-drills:11}
\textbf{Feuille 11.} Corriger la feuille avec cette clé. Votre score est le nombre de bonnes réponses moins trois fois le nombre de mauvaises.
\par\nopagebreak
% Generated by code/interviews-1/03-speed-drills/python/gen_drill.py. Do not edit. French copy by draft/pilot-quant-fr/ivgen_fr.py.
\begin{ivdrillkey}
\ivkey{1}{$91.2$} & \ivkey{6}{$31096$} & \ivkey{11}{$25149$} & \ivkey{16}{$356409$} & \ivkey{21}{$146.52$} \\
\ivkey{2}{$38624$} & \ivkey{7}{$93704$} & \ivkey{12}{$148$} & \ivkey{17}{$66.96$} & \ivkey{22}{$185.25$} \\
\ivkey{3}{$813604$} & \ivkey{8}{$23643$} & \ivkey{13}{$31710$} & \ivkey{18}{$27417$} & \ivkey{23}{$21614$} \\
\ivkey{4}{$9604$} & \ivkey{9}{$334$} & \ivkey{14}{$40804$} & \ivkey{19}{$29406$} & \ivkey{24}{$33222$} \\
\ivkey{5}{$315$} & \ivkey{10}{$51894$} & \ivkey{15}{$219$} & \ivkey{20}{$48654$} &  \\
\end{ivdrillkey}


\textbf{Quand laisser une case vide.} Avec $+1$ et $-3$, une réponse devinée, juste avec la probabilité $p$, ajoute en moyenne $4p - 3$, de sorte qu'elle ne paie que lorsque $p > \tfrac34$. Si vous n'êtes pas assez sûr d'une question, laissez la case vide.

\textbf{Le rythme.} Six minutes pour 24 questions font 15 secondes par question. C'est une feuille de précision. Elle n'est pas faite pour être terminée au détriment du soin, et le travail au brouillon est permis.

\textbf{Où l'on perd des points.} Repérer la famille de chaque question manquée ou passée.
% Generated by code/interviews-1/03-speed-drills/python/gen_drill.py. Do not edit. French copy by draft/pilot-quant-fr/ivgen_fr.py.
\ivfamily{Trois chiffres fois deux chiffres}{2, 8, 10, 11, 18, 19, 20, 24}{Décomposer le facteur à deux chiffres, et additionner avec soin. Puis vérifier le chiffre des unités.}
\ivfamily{Autour des centaines}{6, 7, 13, 23}{$104 \times 297 = 104 \times 300 - 104 \times 3$.}
\ivfamily{Carrés autour des centaines}{3, 4, 14, 16}{$(100k + d)^2 = 10000k^2 + 200kd + d^2$, avec $d$ petit et éventuellement négatif.}
\ivfamily{Division longue}{5, 9, 12, 15}{Retirer des multiples ronds du diviseur, et tenir le décompte du quotient.}
\ivfamily{Produits de décimaux}{1, 17, 21, 22}{Compter les décimales, multiplier comme des entiers, et vérifier l'ordre de grandeur avec une somme arrondie.}


\iqwrong{Se précipiter pour finir. Une série de 24 questions tentées dont 4 fausses donne $20 - 12 = 8$, et une série de 15 tentées sans erreur donne 15. À ce tarif, une seule inattention annule trois bonnes réponses.}
\iqfollow{Comment vérifier une question sans la refaire entièrement~?}{Vérifier le chiffre des unités, puis la somme des chiffres modulo 9. Pour $378 \times 46 = 17\,388$, les sommes des chiffres donnent $0 \times 1 = 0$, et la somme des chiffres du résultat vaut 27, soit 0 modulo 9. Une telle vérification ne manque que l'inversion de deux chiffres ou une erreur égale à un multiple de 9.}
\end{solution}

\begin{solution}{iq:iv1:speed-drills:12}
\textbf{Feuille 12.} Corriger la feuille avec cette clé. Votre score est le nombre de bonnes réponses moins le double du nombre de mauvaises, comme pour la feuille 1, afin que vous puissiez comparer directement les deux scores.
\par\nopagebreak
% Generated by code/interviews-1/03-speed-drills/python/gen_drill.py. Do not edit. French copy by draft/pilot-quant-fr/ivgen_fr.py.
\begin{ivdrillkey}
\ivkey{1}{$4451$} & \ivkey{17}{$9996$} & \ivkey{33}{$4.8$} & \ivkey{49}{$570$} & \ivkey{65}{$14292$} \\
\ivkey{2}{$13$} & \ivkey{18}{$1575$} & \ivkey{34}{$792$} & \ivkey{50}{$2142$} & \ivkey{66}{$3552$} \\
\ivkey{3}{$65$} & \ivkey{19}{$1024$} & \ivkey{35}{$3484$} & \ivkey{51}{$\frac{32}{27}$}\rule[-9.5pt]{0pt}{23pt} & \ivkey{67}{$\frac{31}{50}$}\rule[-9.5pt]{0pt}{23pt} \\
\ivkey{4}{$289$} & \ivkey{20}{$12$} & \ivkey{36}{$6561$} & \ivkey{52}{$4225$} & \ivkey{68}{$6489$} \\
\ivkey{5}{$5$} & \ivkey{21}{$2409$} & \ivkey{37}{$2368$} & \ivkey{53}{$236.5$} & \ivkey{69}{$4$} \\
\ivkey{6}{$38$} & \ivkey{22}{$266.5$} & \ivkey{38}{$1588$} & \ivkey{54}{$125$} & \ivkey{70}{$21$} \\
\ivkey{7}{$1849$} & \ivkey{23}{$9975$} & \ivkey{39}{$7569$} & \ivkey{55}{$256$} & \ivkey{71}{$0.104$} \\
\ivkey{8}{$1148$} & \ivkey{24}{$2730$} & \ivkey{40}{$256$} & \ivkey{56}{$168$} & \ivkey{72}{$2398$} \\
\ivkey{9}{$\frac{7}{8}$}\rule[-9.5pt]{0pt}{23pt} & \ivkey{25}{$3988$} & \ivkey{41}{$99$} & \ivkey{57}{$35$} & \ivkey{73}{$4884$} \\
\ivkey{10}{$0.6875$} & \ivkey{26}{$1320$} & \ivkey{42}{$168$} & \ivkey{58}{$1408$} & \ivkey{74}{$9566$} \\
\ivkey{11}{$0.3125$} & \ivkey{27}{$56$} & \ivkey{43}{$495$} & \ivkey{59}{$12081$} & \ivkey{75}{$286$} \\
\ivkey{12}{$0.275$} & \ivkey{28}{$9984$} & \ivkey{44}{$8723$} & \ivkey{60}{$4990$} & \ivkey{76}{$0.4375$} \\
\ivkey{13}{$392$} & \ivkey{29}{$2952$} & \ivkey{45}{$11394$} & \ivkey{61}{$86$} & \ivkey{77}{$84$} \\
\ivkey{14}{$2916$} & \ivkey{30}{$4200$} & \ivkey{46}{$1850$} & \ivkey{62}{$8$} & \ivkey{78}{$70$} \\
\ivkey{15}{$10.2$} & \ivkey{31}{$16384$} & \ivkey{47}{$2209$} & \ivkey{63}{$3025$} & \ivkey{79}{$0.35$} \\
\ivkey{16}{$3432$} & \ivkey{32}{$19$} & \ivkey{48}{$330$} & \ivkey{64}{$\frac{1}{125}$}\rule[-9.5pt]{0pt}{23pt} & \ivkey{80}{$0.625$} \\
\end{ivdrillkey}


\textbf{Quand laisser une case vide.} Avec $+1$ et $-2$, une réponse devinée, juste avec la probabilité $p$, ajoute en moyenne $3p - 2$, de sorte qu'elle paie lorsque $p > \tfrac23$.

\textbf{Le rythme.} Huit minutes pour 80 questions font 6 secondes par question, comme pour la feuille 1, avec des questions plus difficiles. Peu de lecteurs la termineront, donc choisissez vos questions. Traitez tout de suite chaque famille que vous maîtrisez dès que vous la reconnaissez, et passez le reste au premier tour.

\textbf{Où l'on perd des points.} Repérer la famille de chaque question manquée ou passée.
% Generated by code/interviews-1/03-speed-drills/python/gen_drill.py. Do not edit. French copy by draft/pilot-quant-fr/ivgen_fr.py.
\ivfamily{Additionner et soustraire de grands nombres}{1, 24, 37, 38, 44, 45, 59, 65, 68, 72, 74}{Arrondir un nombre puis corriger.}
\ivfamily{Deux chiffres fois deux chiffres}{8, 16, 18, 21, 26, 29, 30, 35, 46, 50, 58, 66, 73}{Décomposer un facteur, ou utiliser la différence de deux carrés si les facteurs sont proches.}
\ivfamily{Carrés}{7, 14, 39, 47, 52, 63}{$n^2 = (n - d)(n + d) + d^2$.}
\ivfamily{Autour d'un nombre rond}{17, 23, 25, 28, 60}{$(100 + d)(100 - d) = 10000 - d^2$.}
\ivfamily{Pourcentages}{4, 13, 15, 22, 27, 33, 40, 41, 42, 49, 53, 56}{Construire à partir de $10\%$ et de $1\%$. Pour une variation, multiplier par $1 + r$.}
\ivfamily{Fractions et décimaux}{9, 10, 11, 12, 51, 71, 76, 79, 80}{Connaître les huitièmes, seizièmes et quarantièmes. Simplifier avant de multiplier.}
\ivfamily{Conversions}{3, 61, 64, 67}{Ramener une fraction à 100 pour obtenir un pourcentage. Lire un décimal sur une puissance de 10 pour obtenir une fraction.}
\ivfamily{Puissances et racines}{6, 19, 20, 31, 36, 54, 55, 57, 62, 77}{Connaître les puissances de 2 jusqu'à $2^{16}$ et les carrés jusqu'à 29.}
\ivfamily{Coefficients binomiaux}{34, 43, 48, 70, 75, 78}{Simplifier avant de multiplier.}
\ivfamily{Algèbre}{2, 5, 32, 69}{Défaire chaque étape mentalement.}


\iqwrong{Tenter des questions au hasard jusqu'à la fin du temps. Cinquante bonnes réponses sans aucune erreur donnent 50, et 65 bonnes réponses avec 10 erreurs donnent 45. Savoir passer est ici une compétence aussi importante que la vitesse.}
\iqfollow{Qu'est-ce qu'un bon score~?}{Ce livre ne donne aucune norme, car il n'en a pas qu'il puisse justifier. Votre propre historique est la mesure utile~: notez votre score à la date, revenez sur les familles que vous avez manquées, et prenez dans une semaine une feuille générée par \texttt{gen\_drill.py} avec une nouvelle graine.}
\end{solution}
